\section{Conclusion}
\label{sec:conclusion}

This study examined two evidence tasks that a pharmaceutical decision-support system should not collapse: retrieval of citable regulatory passages and verification of structured drug-shortage facts. We built a hierarchical corpus with source-preserving summaries and tables, connected it to a provenance-bearing shortage graph, and evaluated direct, hierarchical, graph, enriched, fused, gated, and neural retrieval variants.

The strongest confirmed source-chunk baseline was context-matched BM25. Graph access did not improve aggregate text ranking, and cumulative flat enrichment reduced dense retrieval performance. The graph did recover six accepted structured paths, which supports its use for entity relations and provenance rather than as a replacement for passage retrieval. Independent double annotation and joint adjudication produced a 58-query Gold set with strong core binary agreement. On that set, BM25 reached Hit@5 of 0.948, MRR of 0.820, and nDCG@5 of 0.800. A locked supplementary Qwen3 reranker increased the point estimates to 0.983, 0.832, and 0.817, but paired uncertainty did not establish a general advantage.

The resulting architecture is deliberately bounded. Lexical retrieval provides the current default ranking, neural and hierarchical components can reorder candidates or expose context after separate validation, and graph paths present structured facts with their sources. The system does not predict shortage duration or causal impact. It provides an auditable interface through which an analyst can inspect what a public record says, which entities it connects, and which regulatory passage supports the accompanying quality-risk interpretation.
